%! AP Statistics — Unit 2, Lesson 2.7 — Exit Ticket — Answer Key
\documentclass[10pt]{article}
\usepackage{apstats-article}
\usepackage{apstats-key}

\begin{document}

\pageheader{Unit 2, Lesson 2.7}{Exit Ticket --- Answer Key}
\namedateperiod

\vspace{0.18in}

% ── SCENARIO ─────────────────────────────────────────────────────────────────
\begin{tcolorbox}[colback=sky, colframe=navy, boxrule=0.9pt, arc=2mm,
    left=4mm, right=4mm, top=2.5mm, bottom=2.5mm]
\small\textit{A nutritionist studies the relationship between daily sugar intake (grams, $x$) and
resting heart rate (bpm, $y$). The regression equation is
$\widehat{\text{HR}} = 58 + 0.2(\text{sugar})$.
One participant consumed \textbf{40 grams} of sugar and had a resting heart rate of \textbf{71 bpm}.}
\end{tcolorbox}

\vspace{0.16in}

\begin{enumerate}[leftmargin=1.6em, itemsep=0.2in]

  \item Calculate the predicted heart rate and the residual for this participant. Show your work.

        \vspace{6pt}
        $\hat{y} = $ \ans{$58 + 0.2(40) = 58 + 8 = 66$ bpm}

        \vspace{6pt}
        $e = y - \hat{y} = $ \ans{$71 - 66 = +5$ bpm}

  \item Interpret the residual in context. State whether the actual value is above or below
        the regression line.\\[8pt]
        \ansline{The actual resting heart rate for this participant (71 bpm) is 5 bpm higher
        than the value predicted by the regression line (66 bpm). The residual is positive,
        so this data point falls above the regression line.}

\end{enumerate}

\vspace{0.16in}

% ── AP-STYLE ─────────────────────────────────────────────────────────────────
\begin{tcolorbox}[colback=redbg, colframe=redacc, boxrule=0.8pt, arc=2mm,
    left=4mm, right=4mm, top=2.5mm, bottom=2.5mm]
\small\textbf{AP-Style Practice}\enspace
A residual plot for a regression of $y$ = fuel efficiency (mpg) on $x$ = engine size (liters)
shows a clear U-shaped pattern. Which conclusion is most appropriate?

\vspace{7pt}
\begin{enumerate}[label=\textbf{(\Alph*)}, leftmargin=2em, itemsep=4pt]
  \item There is no relationship between engine size and fuel efficiency.
  \item A linear model is an appropriate model for the data because residuals are shown.
  \item A linear model may not be an appropriate model; the curved pattern suggests a
        nonlinear model may provide a better fit.
        \textcolor{keyred}{\textbf{$\leftarrow$ correct}}
  \item The regression line has a negative slope.
\end{enumerate}
\end{tcolorbox}

\vspace{0.08in}
\noindent\textcolor{slate}{\small My answer: \ans{C} \quad Because:
\ans{A curved residual plot means the linear model misses systematic structure in the data.
(A) is wrong --- there is a relationship, it's just nonlinear. (B) confuses constructing a
residual plot with the conclusion. (D) is about the slope of the original regression, not
the residual plot.}}

\begin{teachernote}
A residual of $+5$ bpm is worth emphasizing: it means the model underpredicted by 5 bpm for
this individual. Students who write $\hat{y} - y$ (negative) have the subtraction reversed
--- address this at item 1 before they carry the error into item 2.
\end{teachernote}

\end{document}
